\section{Continuum Derivation Details}
\label{app:continuum-derivation-details}

This appendix supplements the continuum mechanics details that support the main-body
assumptions without interrupting the reduced-order forward-model
narrative within this report. It is appendix background for notation and
model hierarchy, not additional evidence for the numerical study.
All identities below are classical identities under the regularity needed for
the displayed derivatives, traces, and changes of variables. In particular,
$\Divbf$ is the row-divergence of a tensor field, and boundary forms assume the
trace and normal hypotheses needed for the displayed integrals. The final
subsection gives compact integral and variational bridge forms used to
distinguish finite-volume, finite-element, and moving-domain discretization
contracts.

\subsection{Material Derivative and Trajectory Details}
\label{subsec:material-derivative-trajectory-details}

For a scalar field $\phi\in C^1(\mathcal{Q}_T)$ and a sufficiently regular flow
map $\La_t$, define the Lagrangian pullback
$\phihat(t,\vxi):=\phi(t,\La_t(\vxi))$. Differentiating with respect to time
along a trajectory and using
$\pd_t \La_t(\vxi)=\vu(t,\La_t(\vxi))$ gives
\begin{flalign*}
    \quad \dt\phi(t,\La_t(\vxi)) = \pd_t \phi(t,\vx)+\nabla\phi(t,\vx)\cdot\vu(t,\vx), \qquad \vx=\La_t(\vxi). &&
\end{flalign*}
The identity
\begin{flalign*}
    \quad \Div(\phi\vu)=\nabla\phi\cdot\vu+\phi\Div\vu &&
\end{flalign*}
also gives
\begin{flalign*}
    \quad D_t\phi=\pd_t \phi+\Div(\phi\vu)-\phi\Div\vu. &&
\end{flalign*}

\subsection{Jacobian Evolution and Reynolds Transport}
\label{subsec:jacobian-evolution-reynolds-transport}

Let $\vFhat_t=\nabla_{\vxi}\La_t$ and $\Jhatdet_t=\det\vFhat_t$. Here $\La_t$
is assumed to be an orientation-preserving $C^1$ change of variables on the
material volume being considered, with the additional time regularity needed to
differentiate $\vFhat_t$ and $\Jhatdet_t$. The standard Jacobi identity
gives
\begin{flalign*}
    \quad
    \pd_t \Jhatdet_t
    =
    \Jhatdet_t\,\mathrm{tr}\!\left(\pd_t \vFhat_t\,\vFhat_t^{-1}\right)
    =
    \Jhatdet_t\,\Div\vu(t,\La_t(\vxi)),
    &&
\end{flalign*}
which is the Lagrangian form of $D_t\Jdet_t=\Jdet_t\Div\vu$
\cite[Part~I, Ch.~1, Sec.~1.3]{GaldiEtAl2008HemodynamicalFlows}.

For $V(t)=\La_t(V(0))$, change variables in
$I(t)=\int_{V(t)}\phi(t,\vx)\dx$:
\begin{flalign*}
    \quad I(t) = \int_{V(0)}\phi(t,\La_t(\vxi))\Jhatdet_t(\vxi)\dxi. &&
\end{flalign*}
Differentiating under the integral sign and substituting the Jacobian-evolution
identity gives
\begin{flalign*}
    \quad
    I'(t)
    =
    \int_{V(t)}
    \left(D_t\phi+\phi\Div\vu\right)\dx
    =
    \int_{V(t)}
    \left(\pd_t \phi+\Div(\phi\vu)\right)\dx. &&
\end{flalign*}
If $V(t)$ has Lipschitz boundary and the trace of $\phi\vu$ is integrable on
$\partial V(t)$, the divergence theorem gives the boundary form
\begin{flalign*}
    \quad
    I'(t)=\int_{V(t)} \pd_t \phi\,\dx
    +\int_{\partial V(t)}\phi\,\vu\cdot\nhat\,dS. &&
\end{flalign*}

\subsection{Linear Momentum Balance Details}
\label{subsec:linear-momentum-balance-details}

Starting from the material-volume balance
\begin{flalign*}
    \quad
    \dt\int_{V(t)}\rho\vu\,dV
    =
    \int_{\partial V(t)}\vT\nhat\,dS
    +\int_{V(t)}\rho\vfb\,dV, &&
\end{flalign*}
Reynolds transport gives
\begin{flalign*}
    \quad
    \dt\int_{V(t)}\rho\vu\,dV
    =
    \int_{V(t)}
    \left(D_t(\rho\vu)+\rho\vu\,\Div\vu\right)dV. &&
\end{flalign*}
Using mass conservation,
$D_t\rho+\rho\Div\vu=0$, this reduces to
\begin{flalign*}
    \quad
    \int_{V(t)}\rho D_t\vu\,dV
    =
    \int_{V(t)}\rho\left(\pd_t \vu+(\vu\cdot\nabla)\vu\right)dV. &&
\end{flalign*}
The surface traction term satisfies
\begin{flalign*}
    \quad
    \int_{\partial V(t)}\vT\nhat\,dS
    =
    \int_{V(t)}\Divbf(\vT)\,dV. &&
\end{flalign*}
This uses the tensor row-divergence convention under the regularity needed for
the stress trace $\vT\nhat$ and the volume integral of $\Divbf\vT$. Since
$V(t)$ is arbitrary within the localization class, the differential form is
\begin{flalign*}
    \quad
    \rho\left(\pd_t \vu+(\vu\cdot\nabla)\vu\right)
    =
    \Divbf(\vT)+\rho\vfb. &&
\end{flalign*}
The conservative Eulerian form
\begin{flalign*}
    \quad
    \pd_t (\rho\vu)+\Divbf(\rho\vu\otimes\vu)
    =
    \Divbf(\vT)+\rho\vfb &&
\end{flalign*}
is equivalent whenever the continuity equation holds, because
\begin{flalign*}
    \quad
    \pd_t (\rho\vu)+\Divbf(\rho\vu\otimes\vu)
    =
    \rho\left(\pd_t \vu+(\vu\cdot\nabla)\vu\right)
    +\vu\left(\pd_t \rho+\Div(\rho\vu)\right). &&
\end{flalign*}

\subsection{Integral and variational bridge forms}
\label{subsec:continuum-integral-variational-bridge}

For a fixed control volume $K$, the finite-volume bridge is the integral
balance
\begin{flalign*}
    \quad
    \frac{d}{dt}\int_K U\dx
    +
    \int_{\partial K}F(U)\cdot\nhat\,dS
    =
    \int_K S(U)\dx, &&
\end{flalign*}
with the numerical flux supplying the discrete approximation to the boundary
integral. This is distinct from a mixed finite-element incompressible weak
form. On a fixed fluid domain $\OmgB$ with boundary parts $\Gamma_D$ and
$\Gamma_N$, the Newtonian schematic weak form is: find $(\vu,p)$ satisfying
the essential velocity data such that
\begin{flalign*}
    \quad
    \begin{aligned}
        \int_{\OmgB}\rho\left(\pd_t\vu+(\vu\cdot\nabla)\vu\right)\cdot\vv\,d\vx
        &+
        \int_{\OmgB}2\eta\vct{D}(\vu):\vct{D}(\vv)\,d\vx \\
        &-
        \int_{\OmgB}p\,\Div\vv\,d\vx
        =
        \int_{\OmgB}\rho\vfb\cdot\vv\,d\vx
        +
        \int_{\Gamma_N}\vt_N\cdot\vv\,dS, \\
        \int_{\OmgB}r\,\Div\vu\,d\vx &=0,
    \end{aligned}
    &&
\end{flalign*}
for admissible velocity tests $\vv$ and pressure tests $r$. The pressure space
must specify a gauge or quotient convention when the boundary conditions leave
pressure determined only up to a constant. In an ALE or moving-domain FSI
variant, the convective derivative is written with relative velocity
$\vu-\vw$, and the fluid--structure interface contributes the kinematic and
traction conditions
\begin{flalign*}
    \quad
    \vu_f=\pd_t\vd_s,
    \qquad
    \vT_f\nhat_f+\vT_s\nhat_s=\zero. &&
\end{flalign*}
These bridge forms identify the mathematical contracts later numerical methods
may discretize; they do not assert that every retained computational workflow
implements the full moving-domain FSI form.
